\subsection{Rappresentazione della mesh mediante un grafo}

La triangolazione $\mathcal T_h$ viene rappresentata mediante un grafo $\mathcal G=(\mathcal V,\mathcal E)$. Ciascun nodo della mesh corrisponde a un vertice del grafo; due vertici sono connessi quando i nodi corrispondenti appartengono allo stesso elemento triangolare. Gli archi sono trattati in entrambe le direzioni. La rappresentazione segue il paradigma delle reti neurali su grafi e dei simulatori graph-based \citep{battaglia2018relational,pfaff2021learning}. La connettività è costruita dalla mesh corrente, mentre le feature geometriche vengono aggiornate a ogni finestra di adattamento.

L'inferenza dei modelli usati nel benchmark PDE riceve sei feature nodali:
\begin{equation}
\bm z_i=
\left[
x_i,\ y_i,\ \widetilde u_i,\ 
\widetilde{\partial_xu}_i,\ 
\widetilde{\partial_yu}_i,\ 
\chi_{\partial\Omega}(i)
\right]^T,
\label{eq:node_features}
\end{equation}
dove $\widetilde u_i$ e le derivate discrete sono le quantità scalate secondo l'adattatore di feature effettivamente impiegato, e $\chi_{\partial\Omega}$ identifica i nodi vincolati sul bordo. Per ogni arco orientato $(i,j)$ sono utilizzate tre feature:
\begin{equation}
\bm e_{ij}=
[x_j-x_i,\ y_j-y_i,\ \norm{\bm x_j-\bm x_i}_2]^T.
\label{eq:edge_features}
\end{equation}
Le derivate del campo sono ricavate dalla soluzione FEM corrente. Non vengono forniti alla rete né la soluzione manufatta esatta, né stati temporali futuri, né gli spostamenti target dell'oracolo, né i parametri fisici $(\sigma,\beta,\bm b,\nu)$ come feature esplicite. Questa scelta definisce il perimetro informativo della GNN e rende il modello dipendente dallo stato numerico corrente, non da informazioni a posteriori.

\subsection{Architettura della rete MeshMover}

La rete segue uno schema \emph{encode--process--decode}. Gli encoder trasformano feature nodali e di arco in rappresentazioni latenti; una successione di sei blocchi di message passing aggrega informazioni dai vicini; il decoder produce, per ciascun nodo, uno spostamento bidimensionale:
\begin{equation}
\widehat{\Delta\bm x}_i
=
\mathcal G_\theta(\mathcal G,\bm Z)_i
\in\R^2.
\label{eq:gnn_output}
\end{equation}
Le configurazioni Stage A e Stage D utilizzano una dimensione latente pari a 128 e MLP a due strati. La rete è condivisa sui nodi e non modifica la connettività della triangolazione. Gli spostamenti dei nodi di bordo sono annullati o vincolati a zero dall'adattatore di inferenza, in modo da mantenere la geometria del bordo.

\subsection{Due strategie di apprendimento}

\paragraph{Stage A: apprendimento supervisionato su campi gaussiani.}
La prima strategia costruisce target di rilocazione minimizzando un funzionale associato all'errore d'interpolazione di campi gaussiani localizzati, con penalità di spostamento e di qualità geometrica. Il modello apprende una mappa dalla mesh e dalle feature del campo agli spostamenti di riferimento. Il dataset contiene 210 campioni: 120 per training, 30 per validation, 30 per test interpolativo e 30 per test estrapolativo.

I target Stage A sono riferimenti numerici a budget finito. La loro accettazione è determinata da criteri di riduzione del funzionale e dell'errore d'interpolazione e dalla validità geometrica.

\paragraph{Stage D: apprendimento trajectory-aware su finestre temporali.}
La seconda strategia considera una finestra temporale di ampiezza $H=0.05$ e apprende una traiettoria di rilocazione associata all'evoluzione del campo manufatto. Il dataset supervisionato comprende 88 campioni: 48 training, 12 validation, 12 test interpolativi e 16 test estrapolativi. I target sono generati con un budget L-BFGS-B limitato a 120 iterazioni.

La loss supervisionata combina errore quadratico sugli spostamenti nodali interni e barriera geometrica:
\begin{align}
\mathcal L_{\mathrm{target}}
&=\frac{1}{N_I}\sum_{i\in\mathcal I}
\norm{\widehat{\Delta\bm x}_i-\Delta\bm x_i^{\,\mathrm{ref}}}_2^2,\\
\mathcal L
&=\lambda_{\mathrm{target}}\mathcal L_{\mathrm{target}}
+\lambda_{\mathrm{barrier}}\mathcal L_{\mathrm{barrier}}.
\label{eq:training_loss}
\end{align}
Nella configurazione Stage D, $\lambda_{\mathrm{target}}=1$ e il peso della barriera è $10^{-4}$; il termine di smoothing è disattivato. La barriera è un regolarizzatore di training e non sostituisce la verifica a runtime.

\subsection{Filtro deterministico di sicurezza geometrica}

La predizione della rete non garantisce, da sola, che la nuova mesh sia ammissibile. Sia $\bm X^n$ la matrice delle coordinate correnti e sia $\Delta\bm X$ lo spostamento proposto. Il filtro considera il percorso lineare
\begin{equation}
\bm X(\alpha)=\bm X^n+\alpha\Delta\bm X,\qquad 0\leq\alpha\leq1,
\label{eq:mesh_path}
\end{equation}
con spostamenti nulli sui nodi di bordo. Per un triangolo orientato $K_e=(a,b,c)$, il Jacobiano geometrico (pari al determinante della matrice dei due lati) è
\begin{equation}
J_e(\alpha)=
\det\!\left[
\bm x_b(\alpha)-\bm x_a(\alpha),\
\bm x_c(\alpha)-\bm x_a(\alpha)
\right].
\end{equation}
Poiché ciascun lato dipende linearmente da $\alpha$, il determinante è un polinomio quadratico:
\begin{equation}
J_e(\alpha)=A_e\alpha^2+B_e\alpha+C_e.
\label{eq:jacobian_quadratic}
\end{equation}
Assumendo che la mesh iniziale soddisfi $J_e(0)>\epsilon$, per ogni elemento si determina il primo attraversamento della soglia $J_e=\epsilon$ lungo il percorso. Il fattore globale è il minimo dei fattori ammissibili:
\begin{equation}
\alpha_{\mathrm{safe}}=\min_{e=1,\ldots,M}\alpha_e^\star,
\qquad
\bm X^{n+1}=\bm X^n+\alpha_{\mathrm{safe}}\Delta\bm X.
\label{eq:safety_alpha}
\end{equation}
La costruzione è deterministica e richiede una scansione degli elementi, quindi costo lineare nel numero di celle. In aritmetica esatta, sotto le ipotesi di triangoli lineari, percorso nodale lineare, mesh iniziale ammissibile e corretta individuazione del primo attraversamento, l'analisi del polinomio controlla il Jacobiano lungo l'intero segmento e non soltanto sul punto finale. La proprietà è condizionata all'ammissibilità della mesh di partenza e alla soglia assoluta adottata.

Nei benchmark sono state osservate minime discrepanze rispetto alla soglia assoluta $\epsilon=10^{-4}$ dell'ordine di $10^{-18}$, compatibili con arrotondamento in doppia precisione; non sono state registrate celle invertite o degeneri dopo il filtro. La soglia, il margine configurato e il criterio per celle degeneri sono mantenuti distinti. Poiché il determinante di un elemento bidimensionale scala con l'area, il valore assoluto di $\epsilon$ deve essere considerato parte della configurazione del benchmark e non una soglia automaticamente trasferibile a mesh di diversa risoluzione. Il filtro non viene rimosso nelle valutazioni di accuratezza.

Nel benchmark qui riportato, tutte le inferenze e le valutazioni partono dalla stessa triangolazione strutturata $32\times32$ e utilizzano coordinate assolute $(x_i,y_i)$ tra le feature nodali. La rappresentazione grafica è una scelta architetturale coerente con la futura estensione a connettività non strutturate.

\begin{proposition}[Ammissibilità lungo il percorso in aritmetica esatta]
Sia $\mathcal T_h$ una triangolazione iniziale con $J_e(0)>\epsilon$ per ogni elemento. Si consideri il percorso \eqref{eq:mesh_path} e si determini, per ciascun elemento, il primo valore $\alpha_e^\star\in(0,1]$ per cui $J_e(\alpha_e^\star)=\epsilon$, ponendo $\alpha_e^\star=1$ se non vi è attraversamento in $[0,1]$. Se il fattore globale è scelto non superiore a $\min_e\alpha_e^\star$ e la ricerca delle radici individua correttamente il primo attraversamento, allora $J_e(\alpha)\geq\epsilon$ per ogni elemento e per ogni $\alpha$ lungo il segmento applicato.
\end{proposition}
